\section{Introduction}

\label{sec:intro}

The emergence of large language models as general-purpose reasoning engines has
catalyzed a new generation of software products that wrap model inference in
task-specific interfaces. Coding assistants like Cursor~\cite{cursor2024} embed
models in editor plugins. Autonomous agents like
Devin~\cite{cognition2024devin} provide browser-based execution environments.
CLI tools like Claude Code~\cite{anthropic2024claude} expose agentic workflows
through a terminal. In each case, the model is treated as an intelligent
component within a pre-existing product architecture.

We argue this framing is architecturally incorrect and quantifiably
suboptimal---not because the products are poorly engineered, but because the
product categories themselves impose structural constraints that prevent
cross-layer optimization. This paper introduces the \textbf{Unified Agent
Harness (UAH)}, a vertically integrated architecture that treats the agent
runtime as the primary abstraction layer, and builds all adjacent concerns---LLM
routing, tool execution, compute provisioning, telemetry, and self-improvement---
as first-class citizens of the same system.

\subsection{The Architectural Mismatch}

Autonomous AI agents are fundamentally \emph{stochastic, open-ended, and
long-horizon} systems. A language model generating a 10,000-step software
engineering trajectory samples from an astronomically large action space,
exhibits non-Markovian behavior through accumulated context, and requires
dynamic resource allocation as task complexity evolves mid-execution.

The software interfaces into which agents are typically embedded---IDE extension
APIs, CI/CD pipeline stages, REST endpoints, chat message schemas---were
designed for \emph{deterministic, bounded, short-horizon} interactions:
single function calls, discrete menu actions, request-response cycles. Forcing
the former into the latter creates what we term the \textbf{architectural
mismatch}: an impedance problem between the stochastic, open-ended semantics of
agent computation and the rigid, bounded contracts of existing software
infrastructure.

The mismatch manifests concretely in four observable failure modes:

\begin{enumerate}[leftmargin=1.5em]
    \item \textbf{Context overflow under bounded interfaces:} Chat APIs impose
          context windows that cause long-horizon tasks to lose critical state.
    \item \textbf{Compute rigidity:} A task that starts as a simple file edit
          may escalate to require a full Linux VM, but the IDE extension has no
          mechanism to provision one mid-session.
    \item \textbf{Siloed observability:} Tool invocations, LLM calls, and
          compute events are logged in separate systems, making end-to-end
          diagnosis impossible.
    \item \textbf{Vendor lock-in prevents optimization:} When LLM routing,
          tool execution, and observability live in different vendor products,
          no single component has visibility into the cross-layer signal needed
          to optimize the joint system.
\end{enumerate}

\subsection{The Vertical Integration Imperative}

The software industry has repeatedly learned that vertical integration enables
performance that horizontal composition cannot match. Apple's silicon-to-OS-to-
application stack consistently outperforms equivalent x86 configurations.
Amazon Web Services, by owning networking, storage, compute, and managed
services in a single stack, can optimize network topology, caching policies, and
VM placement globally. The performance gap widens as the number of interacting
components grows.

For autonomous agents, the relevant components span seven distinct layers:
LLM access and routing, tool protocol and execution, agent runtime and
orchestration, channel integration and identity, compute tier management,
self-improvement loops, and observability infrastructure. No existing product
controls all seven. The UAH does.

\subsection{Contributions}

This paper makes the following technical contributions:

\begin{enumerate}[leftmargin=1.5em]
    \item \textbf{Seven-Layer Architecture (Section~\ref{sec:architecture}):}
          A formal specification of the seven layers required for production
          autonomous agents, each defined as a typed interface with input/output
          signatures and invariant conditions.
    \item \textbf{Tiered Compute Runtime (Section~\ref{sec:tiers}):}
          A five-tier compute hierarchy from shared gateway to local desktop,
          with formal state-machine transitions and context migration protocols.
    \item \textbf{Cross-Layer Optimization (Section~\ref{sec:optimization}):}
          Formal proof that joint optimization across layers strictly dominates
          independent per-layer optimization, with empirical validation.
    \item \textbf{Competitive Analysis (Section~\ref{sec:competitive}):}
          A systematic comparison against six major agent platforms across all
          seven layers.
    \item \textbf{Multi-Domain Generalization (Section~\ref{sec:domains}):}
          Demonstration that the seven-layer harness is domain-agnostic,
          applying unchanged to software engineering, marketing, DevOps, and
          research agents.
    \item \textbf{Production Implementation (Section~\ref{sec:implementation}):}
          A description of the Hanzo UAH deployment on Kubernetes with NATS,
          Temporal, Qdrant, and a full identity and key management stack.
\end{enumerate}

\paragraph{Paper Outline.}
Section~\ref{sec:related} reviews related work. Section~\ref{sec:architecture}
defines the seven-layer model. Section~\ref{sec:tiers} formalizes the tiered
compute runtime. Section~\ref{sec:optimization} proves and demonstrates
cross-layer optimization. Section~\ref{sec:competitive} provides the competitive
landscape. Section~\ref{sec:domains} shows multi-domain generalization.
Section~\ref{sec:implementation} describes the production implementation.
Section~\ref{sec:discussion} discusses future directions.
Section~\ref{sec:conclusion} concludes.

% ============================================================================
% 2. RELATED WORK
% ============================================================================
